% ============================================================
%  ch5 Implementation and Optimization.tex
%  Chapter 5: Implementation and Optimization
%  Master's Thesis -- SQL MATCH_RECOGNIZE on Pandas DataFrames
% ============================================================
% Provide \mr if not already defined by the main document
\providecommand{\mr}{\texttt{MATCH\_RECOGNIZE}}
\chapter{Implementation and Optimization}
\label{chap:implementation}

This chapter describes how the logical design introduced in
Chapter~\ref{chap:methodology} is realised as an executable
\texttt{MATCH\_RECOGNIZE} engine for Pandas DataFrames.  The focus is
on the implementation layers that translate a validated logical plan
into executable pattern matching: automata compilation, pattern
tokenisation, NFA/DFA construction, Pandas-based execution, measure
materialisation, and performance optimisation.

The chapter is organised as follows.  Section~\ref{sec:meth-automata}
explains how row patterns are compiled into finite automata.
Section~\ref{sec:meth-pandas} describes execution over Pandas
DataFrames, including partitioning, row classification, skip modes,
output modes, and measure evaluation.  Section~\ref{sec:meth-impl}
summarises the concrete implementation components.  Finally,
Section~\ref{sec:bg-performance} presents the optimisation layer,
including LRU pattern caching, memory management, and vectorised
condition evaluation.

% ----------------------------------------------------------------
\section{Pattern Compilation with Finite Automata}
\label{sec:meth-automata}
% ----------------------------------------------------------------

Pattern matching in \texttt{MATCH\_RECOGNIZE} is fundamentally a
problem of sequence recognition.  A sequence of data rows, each
labelled with the pattern variable(s) it satisfies, constitutes a
string over the alphabet of pattern variables.  The problem of
determining whether this string is accepted by the pattern is
exactly the problem solved by finite automata~\cite{HopcroftUllman}.

% ................................
\subsection{Non-deterministic Finite Automata (NFA)}
\label{subsec:meth-nfa}
% ................................

A \emph{Non-deterministic Finite Automaton} (NFA) is formally
defined as a 5-tuple $M = (Q,\,\Sigma,\,\delta,\,q_0,\,F)$, where:

\begin{itemize}
  \item $Q$ is a finite set of \emph{states};
  \item $\Sigma$ is the \emph{input alphabet} (pattern variables in
    this setting);
  \item $\delta : Q \times (\Sigma \cup \{\varepsilon\}) \to 2^{Q}$
    is the \emph{transition function}, mapping each state and symbol
    (or the empty string $\varepsilon$) to a \emph{set} of successor
    states;
  \item $q_0 \in Q$ is the \emph{initial state}; and
  \item $F \subseteq Q$ is the set of \emph{accepting states}.
\end{itemize}

NFAs naturally model alternation and optional constructs owing to
their ability to be in multiple states simultaneously and to use
$\varepsilon$-transitions (transitions that consume no input).

Figure~\ref{fig:meth-nfa} shows the NFA for pattern $A^{+}\,B?$.

\begin{figure}[htbp]
  \centering
  \begin{tikzpicture}[
      >=stealth, shorten >=1pt,
      node distance=2.8cm, auto, thick,
      every state/.style={circle, draw, minimum size=1cm}
    ]
    \node[state, initial, initial text={}]  (q0) {$q_0$};
    \node[state]                            (q1) [right of=q0] {$q_1$};
    \node[state]                            (q2) [right of=q1] {$q_2$};
    \node[state, accepting]                 (q3) [right of=q2] {$q_3$};

    \path[->]
      (q0) edge              node {$A$}           (q1)
      (q1) edge [loop above] node {$A$}           (q1)
      (q1) edge              node {$\varepsilon$} (q2)
      (q2) edge              node {$B$}           (q3)
      (q2) edge [bend right=35] node [below] {$\varepsilon$} (q3);
  \end{tikzpicture}
  \caption{NFA for the pattern $A^{+}\,B?$ illustrating
    $\varepsilon$-transitions and non-determinism}
  \label{fig:meth-nfa}
\end{figure}

% ................................
\subsection{Thompson's Construction}
\label{subsec:meth-thompson}
% ................................

Thompson's construction algorithm~\cite{Thompson1968} converts a
regular expression into an NFA in linear time $O(n)$, where $n$ is
the length of the expression.  The algorithm proceeds recursively
over the expression structure according to the rules in
Algorithm~\ref{alg:thompson}.

\begin{algorithm}[htbp]
\caption{Thompson's NFA Construction}
\label{alg:thompson}
\begin{algorithmic}[1]
\Procedure{Build}{$e$}
  \If{$e = \varepsilon$}
    \State \Return NFA with one $\varepsilon$-transition from start to accept
  \ElsIf{$e = a \in \Sigma$}
    \State \Return NFA with one $a$-labelled transition from start to accept
  \ElsIf{$e = r \cdot s$ \Comment{concatenation}}
    \State $N_r \gets \Call{Build}{r}$
    \State $N_s \gets \Call{Build}{s}$
    \State merge accept state of $N_r$ with start state of $N_s$
    \State \Return combined NFA
  \ElsIf{$e = r \mid s$ \Comment{alternation}}
    \State $N_r \gets \Call{Build}{r}$
    \State $N_s \gets \Call{Build}{s}$
    \State add new start $q_\text{new}$ with $\varepsilon$-edges to $q_0^{N_r}$ and $q_0^{N_s}$
    \State add new accept $f_\text{new}$ with $\varepsilon$-edges from $F^{N_r}$ and $F^{N_s}$
    \State \Return combined NFA
  \ElsIf{$e = r^{*}$ \Comment{Kleene star}}
    \State $N_r \gets \Call{Build}{r}$
    \State add $\varepsilon$-edge: start $\to$ accept \Comment{zero repetitions}
    \State add $\varepsilon$-edge: $F^{N_r} \to q_0^{N_r}$ \Comment{loop back}
    \State \Return modified NFA
  \ElsIf{$e = r^{\{n,m\}}$ \Comment{bounded quantifier}}
    \State construct $n$ mandatory copies followed by $(m{-}n)$ optional copies
    \State \Return combined NFA
  \EndIf
\EndProcedure
\end{algorithmic}
\end{algorithm}

The resulting NFA has at most $2n$ states and is constructed in
$O(n)$ time and space, making it suitable for patterns of arbitrary
complexity~\cite{DragonBook}.

% ................................
\subsection{Pattern Tokenisation}
\label{subsec:meth-pattern-tokeniser}
% ................................

Before an NFA can be constructed the raw pattern string extracted
from the SQL query must be converted into a structured, traversable
representation.  The \texttt{tokenize\_pattern()} function in
\texttt{src/matcher/pattern\_tokenizer.py} performs this step in a
single left-to-right scan, producing a flat list of
\texttt{PatternToken} objects—the \emph{linearised} form of the
pattern's syntax tree.

\paragraph{PatternToken structure.}
Each \texttt{PatternToken} is a Python \texttt{dataclass} with the
following key fields:

\begin{itemize}
  \item \textbf{\texttt{type}} — a \texttt{PatternTokenType} enum
    value determining which NFA fragment builder is invoked
    (see Table~\ref{tab:token-types});
  \item \textbf{\texttt{value}} — the raw textual value extracted
    from the pattern string (variable name, operator character, etc.);
  \item \textbf{\texttt{quantifier}} — optional quantifier string
    (\texttt{'*'}, \texttt{'+'}, \texttt{'?'}, or
    \texttt{'\{$n$,$m$\}'}) copied from the immediately following
    quantifier when present; defaults to \texttt{None};
  \item \textbf{\texttt{greedy}} — boolean flag indicating greedy
    (\texttt{True}, the default) or reluctant (\texttt{False})
    quantification, set when a trailing \texttt{?} follows the
    quantifier;
  \item \textbf{\texttt{metadata}} — a dictionary carrying
    type-specific structured data (e.g.\ a \texttt{variables} list
    for \texttt{PERMUTE} tokens, an \texttt{excluded\_variables}
    list for \texttt{EXCLUSION} tokens); and
  \item \textbf{\texttt{priority}} — an integer priority inherited
    from the token's position in the stream, propagated to NFA
    transitions to enforce greedy/leftmost semantics.
\end{itemize}

\begin{table}[htbp]
  \centering
  \caption{The thirteen \texttt{PatternTokenType} values and their
    roles in the pattern string}
  \label{tab:token-types}
  \begin{tabular}{lp{8cm}}
    \toprule
    \textbf{Token type}         & \textbf{Role in pattern} \\
    \midrule
    \texttt{LITERAL}            & A pattern variable name (\texttt{A}, \texttt{UP},
                                  \texttt{STRT}, etc.) \\
    \texttt{ALTERNATION}        & The \texttt{|} operator \\
    \texttt{PERMUTE}            & A complete \texttt{PERMUTE(\ldots)} construct
                                  carrying its variable list in \texttt{metadata} \\
    \texttt{GROUP\_START}       & An opening parenthesis \texttt{(} \\
    \texttt{GROUP\_END}         & A closing parenthesis \texttt{)} \\
    \texttt{ANCHOR\_START}      & The start-of-partition anchor \texttt{\^{}} \\
    \texttt{ANCHOR\_END}        & The end-of-partition anchor \texttt{\$} \\
    \texttt{EXCLUSION}          & A complete \texttt{\{- \ldots -\}} exclusion block \\
    \texttt{EXCLUSION\_START}   & Opening delimiter \texttt{\{-} \\
    \texttt{EXCLUSION\_END}     & Closing delimiter \texttt{-\}} \\
    \texttt{QUANTIFIER}         & A standalone quantifier node \\
    \texttt{SUBSET}             & A \texttt{SUBSET} variable reference \\
    \texttt{SEQUENCE}           & A composite sequence node \\
    \bottomrule
  \end{tabular}
\end{table}

\paragraph{Tokenisation algorithm.}
The internal function \texttt{\_tokenize\_pattern\_internal()} scans
the pattern string character by character, dispatching on the current
position:

\begin{enumerate}
  \item Whitespace is skipped unconditionally.
  \item The keyword \texttt{PERMUTE} (case-insensitive) triggers
    \texttt{\_parse\_permute\_pattern()}, which collects the
    comma-separated variable list inside the following parentheses
    and returns a single \texttt{PERMUTE} token whose
    \texttt{metadata[`variables']} holds the parsed variable tokens.
  \item The two-character sequence \texttt{\{-} triggers
    \texttt{\_parse\_exclusion\_pattern()}, which reads to the
    matching \texttt{-\}} (tracking nesting depth) and returns an
    \texttt{EXCLUSION} token.
  \item The structural characters \texttt{(}, \texttt{)},
    \texttt{|}, \texttt{\^{}}, and \texttt{\$} each produce a
    single-character token of the corresponding type.
  \item Quantifier characters \texttt{*}, \texttt{+}, \texttt{?},
    and \texttt{\{} are parsed by \texttt{parse\_quantifier\_at()}
    and \emph{attached} to the preceding token as its
    \texttt{quantifier} and \texttt{greedy} fields—no separate
    token is inserted.
  \item Identifier characters (\texttt{[A-Za-z0-9\_]}) are
    accumulated into a \texttt{LITERAL} token.
\end{enumerate}

After the scan, \texttt{\_optimize\_token\_sequence()} makes a
conservative pass that removes structural redundancies while
preserving the one-variable-per-token invariant that the NFA builder
requires: adjacent \texttt{LITERAL} tokens that look like pattern
variable names are deliberately \emph{not} merged.

\paragraph{Quantifier normalisation.}
Quantifiers are stored as attributes of the token they modify rather
than as separate tokens.  The helper \texttt{parse\_quantifier()}
normalises all surface forms to the triple
$(\textit{min},\,\textit{max},\,\textit{greedy})$:
\begin{center}
\begin{tabular}{lll}
  \texttt{?} & $\to$ & $(0,\;1,\;\mathit{True})$ \\
  \texttt{*} & $\to$ & $(0,\;\infty,\;\mathit{True})$ \\
  \texttt{+} & $\to$ & $(1,\;\infty,\;\mathit{True})$ \\
  \texttt{\{$n$\}} & $\to$ & $(n,\;n,\;\mathit{True})$ \\
  \texttt{\{$n$,$m$\}} & $\to$ & $(n,\;m,\;\mathit{True})$ \\
  append \texttt{?} to any & $\to$ & $\mathit{greedy} = \mathit{False}$
\end{tabular}
\end{center}

\noindent This triple is consumed directly by
\texttt{NFABuilder.\_apply\_quantifier()} during NFA construction.

% ................................
\subsection{NFA Construction from Tokens}
\label{subsec:meth-nfa-from-tokens}
% ................................

The \texttt{NFABuilder} class in \texttt{src/matcher/automata.py}
implements the \emph{visitor} pattern over the token list.  Its
central method \texttt{\_process\_sequence()} iterates the token
stream and dispatches on each token's \texttt{type} field to invoke
the appropriate NFA-fragment builder—mirroring the recursive
tree-traversal structure of classical Thompson's construction but
applied to the flat linearised token representation rather than an
explicit parse tree.

\paragraph{Overall build sequence.}
The public entry point \texttt{NFABuilder.build()} executes the
following steps:

\begin{enumerate}
  \item Check the LRU pattern cache (keyed on an MD5 hash of the
    token list and DEFINE context) for a previously compiled NFA.
    A cache hit returns a deep clone of the stored NFA to prevent
    state sharing.
  \item On a cache miss call \texttt{\_build\_nfa\_internal()}, which:
    \begin{enumerate}[label=(\alph*)]
      \item validates anchor semantics (e.g.\ rejecting patterns such
        as \texttt{A\^{}} where a start anchor follows a literal);
      \item creates a global start state $q_\text{start}$ and accept
        state $q_\text{accept}$;
      \item delegates to \texttt{\_process\_sequence()} for the
        pattern body; and
      \item connects the pattern fragment to the global states via
        $\varepsilon$-transitions.
    \end{enumerate}
  \item Apply \texttt{NFA.optimize()} to remove unreachable states,
    deduplicate transitions, and pre-compute DFA hot-paths.
  \item Cache the compiled NFA and return it.
\end{enumerate}

\subsubsection{Token-to-NFA Fragment Mapping}
\label{subsubsec:meth-token-to-nfa}

Table~\ref{tab:token-nfa-map} summarises how the visitor translates
each token type into NFA structure.

\begin{table}[htbp]
  \centering
  \caption{Token types and their NFA fragment builders in
    \texttt{NFABuilder.\_process\_sequence()}}
  \label{tab:token-nfa-map}
  \begin{tabular}{lp{8.2cm}}
    \toprule
    \textbf{Token type}            & \textbf{NFA construction rule} \\
    \midrule
    \texttt{LITERAL}               & \texttt{create\_var\_states()} produces a
                                     two-state fragment
                                     $q_s \xrightarrow{C(v)} q_e$, where
                                     $C(v)$ is the callable compiled from the
                                     \texttt{DEFINE} predicate for variable $v$;
                                     the token's \texttt{quantifier} (if any) is
                                     applied by \texttt{\_apply\_quantifier()} \\[4pt]
    \texttt{ALTERNATION}           & The left branch is finalised; the right branch
                                     is processed by a recursive call to
                                     \texttt{\_process\_sequence()}; a fresh start
                                     state fans out with $\varepsilon$-transitions
                                     to both branch heads, and a fresh end state
                                     collects both branch tails \\[4pt]
    \texttt{GROUP\_START}          & Triggers a recursive call to
                                     \texttt{\_process\_sequence()} that reads
                                     until the matching \texttt{GROUP\_END}; the
                                     group quantifier (if any) is applied to the
                                     resulting fragment \\[4pt]
    \texttt{PERMUTE}               & \texttt{\_process\_permute()} expands all
                                     $n!$ orderings of the $n$ variables and
                                     builds a union of $n!$ concatenation
                                     fragments under a shared start/end pair \\[4pt]
    \texttt{EXCLUSION}             & \texttt{\_process\_exclusion\_token()} builds
                                     the contained sub-pattern normally but marks
                                     every state in the range with
                                     \texttt{is\_excluded = True}; the matcher
                                     filters excluded variables from output at
                                     match time \\[4pt]
    \texttt{ANCHOR\_START}/\texttt{END}
                                   & Creates a single anchor state with
                                     \texttt{is\_anchor = True} and the
                                     corresponding \texttt{anchor\_type};
                                     partition-boundary semantics are enforced
                                     by the matcher \\
    \bottomrule
  \end{tabular}
\end{table}

\subsubsection{Quantifier Expansion}
\label{subsubsec:meth-apply-quantifier}

\texttt{\_apply\_quantifier(start, end, }$\mathit{min}$\texttt{, }$\mathit{max}$\texttt{, }$\mathit{greedy}$\texttt{)}
wraps a fragment $(s,\,e)$ in outer envelope states $(q_s,\,q_e)$
and adds $\varepsilon$-transitions accordingly:

\begin{itemize}
  \item \textbf{\texttt{?}} $(\min=0,\,\max=1)$: adds
    $\varepsilon(q_s \to s)$ for the match case and
    $\varepsilon(q_s \to q_e)$ for the skip case.
  \item \textbf{\texttt{*}} $(\min=0,\,\max=\infty)$: adds
    $\varepsilon(q_s \to s)$, $\varepsilon(e \to q_e)$ (exit),
    $\varepsilon(q_s \to q_e)$ (skip), and
    $\varepsilon(e \to s)$ (loop for repetition).
  \item \textbf{\texttt{+}} $(\min=1,\,\max=\infty)$: as \texttt{*}
    but without the skip transition, enforcing at least one match.
  \item \textbf{\texttt{\{$n$,$m$\}}}: constructs $n$ mandatory
    copies of the fragment in series (connected by
    $\varepsilon$-transitions), followed by $(m-n)$ optional copies
    each of which carries an $\varepsilon$-bypass to $q_e$—the
    standard elimination-of-quantifier reduction.
  \item \textbf{Reluctant} (\texttt{greedy = False}): the same
    $\varepsilon$ structure is used, but the loop-back and
    extension transitions are annotated with a \emph{higher}
    numeric priority value so the epsilon-closure traversal
    prefers the exit path, yielding shorter matches.
\end{itemize}

\subsubsection{PERMUTE Expansion}
\label{subsubsec:meth-permute-expansion}

The \texttt{PERMUTE}$(v_1, v_2, \ldots, v_n)$ construct is
semantically equivalent to the alternation of all $n!$ orderings of
its arguments.  \texttt{\_process\_permute()} proceeds in three
phases:

\begin{enumerate}
  \item \textbf{Complexity guard.}  Patterns with more than eight
    variables are rejected immediately (raising
    \texttt{RuntimeError}) because $8! = 40\,320$ is the largest
    feasible inline expansion; patterns estimated to produce more
    than 50\,000 combinations are redirected to a conservative
    fallback via \texttt{\_process\_permute\_limited()}.
  \item \textbf{Permutation generation.}
    \texttt{ProductionPermuteHandler.expand\_permutation()} is called,
    which uses \texttt{itertools.per\-muta\-tions} for $n \leq 6$ and
    a memory-efficient streaming algorithm for larger $n$.  Results
    are memoised in a thread-safe LRU cache keyed on the sorted
    variable tuple.
  \item \textbf{NFA assembly.}  For each permutation
    $\sigma = (\sigma_1, \ldots, \sigma_n)$, the builder constructs
    a concatenation fragment by calling
    \texttt{create\_var\_states()} for each $\sigma_i$ and connecting
    the fragments with $\varepsilon$-transitions.  All $n!$
    concatenation fragments are then united under a common fork/merge
    pair with the same $\varepsilon$-alternation structure used for
    the \texttt{|} operator.
\end{enumerate}

% ................................
\subsection{Properties of the Constructed NFA}
\label{subsec:meth-nfa-properties}
% ................................

The NFA returned by \texttt{NFABuilder.build()} satisfies the
following formal properties:

\begin{enumerate}
  \item \textbf{Linear size.}  At most two states per
    \texttt{LITERAL} token (the transition's source and target)
    plus a bounded constant of intermediate states per structural
    token (alternation fork/merge, quantifier envelope, etc.).
    Total state count is $O(n)$ in token-list length $n$.

  \item \textbf{Single accept state.}  Exactly one accept state
    $q_\text{accept}$ exists; all fragment tails are connected to
    it via $\varepsilon$-transitions.

  \item \textbf{Priority-annotated $\varepsilon$-transitions.}
    Every $\varepsilon$-transition carries an integer priority
    stored in the \texttt{epsilon\_priorities} dictionary on its
    source \texttt{NFAState}.  Lower numeric values denote higher
    priority.  For greedy quantifiers the loop-back
    $\varepsilon$ has lower priority than the exit $\varepsilon$,
    so the $\varepsilon$-closure computation (which sorts targets by
    ascending priority) extends the match before terminating.
    Reluctant quantifiers reverse the ordering.  These priorities
    are propagated through subset construction so that each DFA
    transition inherits the minimum (highest-priority) value of the
    contributing NFA transitions, implementing SQL:2016 greedy and
    leftmost-match semantics without modifying the DFA construction
    algorithm.

  \item \textbf{Thread safety.}  Each \texttt{NFAState} and the
    \texttt{NFA} object itself are guarded by a
    \texttt{threading.RLock()}.  Epsilon-closure results are cached
    in a bounded dictionary (capped at 1\,000 entries) accessed
    under the same lock, enabling safe concurrent queries by
    multiple partition workers.

  \item \textbf{Structural validation.}  \texttt{NFA.validate()}
    verifies that (i)~all transition targets are valid state
    indices, (ii)~the accept state is reachable from the start
    state, and (iii)~PERMUTE metadata is internally consistent
    before the NFA is returned to the caller.
\end{enumerate}

% ................................
\subsection{Worked Example: Pattern \texttt{A\;B+\;C?}}
\label{subsec:meth-nfa-example}
% ................................

Consider the pattern \texttt{A B+ C?} with three pattern variables
\texttt{A}, \texttt{B}, \texttt{C} each associated with an
arbitrary Boolean DEFINE predicate.  The tokeniser produces the
flat list:
\begin{center}
  \texttt{[\,LITERAL('A'),\quad LITERAL('B', quantifier='+'),\quad
    LITERAL('C', quantifier='?')\,]}
\end{center}
The \texttt{NFABuilder} then processes this list as follows.

\begin{enumerate}
  \item A global start state $q_0$ and accept state $q_a$ are
    created; $q_a$ is marked \texttt{is\_accept = True}.

  \item \textbf{LITERAL \texttt{A} (no quantifier).}
    \texttt{create\_var\_states} allocates $q_1$ and $q_2$; adds
    transition $q_1 \xrightarrow{C(A)} q_2$.
    The visitor adds $\varepsilon(q_0 \to q_1)$ and advances
    \texttt{current} to $q_2$.

  \item \textbf{LITERAL \texttt{B} with quantifier \texttt{+}.}
    \texttt{create\_var\_states} allocates $q_3, q_4$ with
    $q_3 \xrightarrow{C(B)} q_4$.
    \texttt{\_apply\_quantifier}$(q_3,\,q_4,\,1,\,\infty,\,\mathit{True})$
    allocates envelope states $q_5, q_6$ and adds:
    $\varepsilon(q_5 \to q_3)$,\;
    $\varepsilon(q_4 \to q_6)$ (exit),\;
    $\varepsilon(q_4 \to q_3)$ (loop, greedy).
    The visitor adds $\varepsilon(q_2 \to q_5)$ and advances to $q_6$.

  \item \textbf{LITERAL \texttt{C} with quantifier \texttt{?}.}
    \texttt{create\_var\_states} allocates $q_7, q_8$ with
    $q_7 \xrightarrow{C(C)} q_8$.
    \texttt{\_apply\_quantifier}$(q_7,\,q_8,\,0,\,1,\,\mathit{True})$
    allocates $q_9, q_{10}$ and adds:
    $\varepsilon(q_9 \to q_7)$ (take $C$),\;
    $\varepsilon(q_8 \to q_{10})$ (exit after $C$),\;
    $\varepsilon(q_9 \to q_{10})$ (skip $C$, bypass).
    The visitor adds $\varepsilon(q_6 \to q_9)$ and advances to $q_{10}$.

  \item $\varepsilon(q_{10} \to q_a)$ connects the pattern's tail
    to the global accept state.
\end{enumerate}

\noindent
The resulting NFA (11 states: $q_0$–$q_{10}$ plus the accept state
shared with $q_{10}$ in this abbreviated numbering) accepts exactly
the language $A \cdot B^{+} \cdot C^{?}$: one row satisfying $A$,
one or more rows satisfying $B$ (the loop $q_4 \to q_3$ implements
the \texttt{+} greedily), and zero or one row satisfying $C$ (the
bypass $q_9 \to q_{10}$ implements the \texttt{?}).

After construction the NFA is passed to \texttt{DFABuilder.build()}
for subset construction as described in the next section.

% ................................
\subsection{Deterministic Finite Automata (DFA)}
\label{subsec:meth-dfa}
% ................................

A \emph{Deterministic Finite Automaton} (DFA) is a special case of
an NFA in which:
\begin{enumerate*}[label=(\roman*)]
  \item there are no $\varepsilon$-transitions, and
  \item for each state $q$ and symbol $a$, $|\delta(q,a)| = 1$
    (exactly one successor state).
\end{enumerate*}
DFAs require $O(n)$ time to process an input of length $n$ with no
backtracking, making them ideal for row-by-row matching.

\paragraph{Subset Construction (NFA $\to$ DFA).}
The \emph{subset construction} (powerset construction) algorithm
converts an NFA to an equivalent DFA.  Each DFA state corresponds
to a \emph{set} of NFA states.  For a DFA state $S$ and input
symbol $a$:

\[
  \delta_{\text{DFA}}(S,\,a)
  \;=\;
  \varepsilon\text{-closure}\!\left(
    \bigcup_{q \in S} \delta_{\text{NFA}}(q,\,a)
  \right)
\]

where $\varepsilon\text{-closure}(T)$ is the set of all NFA states
reachable from any state in $T$ by zero or more
$\varepsilon$-transitions.  Although the resulting DFA may have up
to $2^{|Q|}$ states in the worst case, patterns arising from
\texttt{MATCH\_RECOGNIZE} syntax are structurally constrained and in
practice yield DFAs with far fewer states.

\paragraph{DFA Minimisation.}
Hopcroft's algorithm~\cite{HopcroftUllman} reduces a DFA to its
canonical minimal form in $O(n \log n)$ time by iteratively
partitioning states into equivalence classes of
indistinguishable states.  The implementation applies minimisation
after subset construction to reduce memory footprint and improve
matching speed.

\paragraph{State Prioritisation for SQL Semantics.}
SQL:2016 mandates \emph{greedy} matching by default: the engine
must prefer the longest possible match, and in the case of
alternation, the leftmost alternative takes precedence.  Standard
DFA theory does not encode these preferences.  The implementation
addresses this by:
\begin{enumerate}
  \item Assigning a \emph{priority} to each NFA transition based on
    its position in the pattern and quantifier greediness.
  \item Propagating priorities through subset construction so that
    each DFA transition inherits the minimum (highest-priority)
    value of the contributing NFA transitions.
  \item Ordering outgoing DFA transitions by priority during
    match execution so that higher-priority paths are explored
    first.
\end{enumerate}


% ----------------------------------------------------------------
\section{Pandas DataFrame Execution Model}
\label{sec:meth-pandas}
% ----------------------------------------------------------------

Pandas~\cite{McKinney2010} is a Python library providing the
\texttt{DataFrame} data structure: a two-dimensional, labelled
table backed by typed NumPy arrays.  Its in-memory, columnar
storage enables vectorised operations over millions of rows, making
it well suited for the row-by-row DFA simulation that
\texttt{MATCH\_RECOGNIZE} requires.

This section describes how the engine transforms a validated
\texttt{MATCH\_RECOGNIZE} clause and a compiled DFA into concrete
output rows.  The process passes through six distinct stages:
partition, order, row classification, DFA simulation,
measure evaluation, and output assembly.
Figure~\ref{fig:exec-pipeline} places these stages in context
of the full five-phase architecture introduced at the start of
this chapter.

\begin{figure}[htbp]
  \centering
  \begin{tikzpicture}[
    node distance=0.6cm and 1.6cm,
    every node/.style={font=\small},
    box/.style={rectangle, draw, rounded corners=3pt,
                minimum width=2.8cm, minimum height=0.75cm,
                align=center},
    arrow/.style={->, thick}
  ]
    \node[box, fill=gray!15] (parse) {SQL Parsing};
    \node[box, fill=gray!15, right=of parse] (sem) {Semantic\\Validation};
    \node[box, fill=gray!15, right=of sem] (plan) {Logical Plan};
    \node[box, fill=gray!15, right=of plan] (comp) {Pattern\\Compilation};
    \node[box, fill=blue!15, right=of comp] (exec) {\textbf{Execution}};
    \draw[arrow] (parse) -- (sem);
    \draw[arrow] (sem) -- (plan);
    \draw[arrow] (plan) -- (comp);
    \draw[arrow] (comp) -- (exec);
  \end{tikzpicture}
  \caption{Five-phase execution pipeline; the execution stage (shaded)
    is the subject of this section}
  \label{fig:exec-pipeline}
\end{figure}

% ................................
\subsection{Partition and Order}
\label{subsec:meth-exec-partition}
% ................................

Execution begins by splitting the input DataFrame $D$ into
independent \emph{partitions}.  The columns named in the
\texttt{PARTITION BY} clause are passed to
\texttt{DataFrame.groupby()}, which returns one sub-frame per
distinct combination of partition-key values.  When no
\texttt{PARTITION BY} clause is present the entire DataFrame
constitutes a single partition.

Each partition is then sorted by the \texttt{ORDER BY} columns
using \texttt{sort\_values()}, establishing the temporal order
within which the DFA simulation will proceed.  The underlying
NumPy index is preserved so that row references in navigation
functions remain stable.

Because partitions share no state, they can be processed
independently—and optionally in parallel (see
Section~\ref{subsec:bg-parallel}).  The remainder of this section
describes the processing of a single partition $P$ of $n$ rows,
indexed $r_0, r_1, \ldots, r_{n-1}$ in sorted order.

% ................................
\subsection{Row Classification}
\label{subsec:meth-exec-classify}
% ................................

Before the DFA simulation begins, every \texttt{DEFINE} predicate
is compiled into a callable Python function.  The
\texttt{compile\_condition()} function in
\texttt{src/matcher/condition\_evaluator.py} parses the predicate
expression—which may reference pattern variables via dot notation
(\texttt{A.price}), navigation functions, and SQL comparison
operators—and wraps it in a closure that accepts a row dictionary
and a \texttt{RowContext} object.

The \texttt{RowContext}
(Listing~\ref{lst:rowcontext-iface}) is the central evaluation
context.  It provides:
\begin{itemize}
  \item \texttt{rows}: the full sorted partition as a list of row
    dictionaries, enabling look-ahead and look-behind via index
    arithmetic;
  \item \texttt{variables}: the current match's variable bindings
    (a mapping from variable name to a list of matching row
    indices);
  \item \texttt{current\_idx}: the index of the row being tested;
  \item \texttt{current\_var}: the name of the pattern variable
    whose predicate is being evaluated; and
  \item \texttt{subsets}: the \texttt{SUBSET} variable definitions.
\end{itemize}

\begin{lstlisting}[
  language=Python,
  caption={Key fields of \texttt{RowContext} consulted during
    condition evaluation},
  label={lst:rowcontext-iface}]
@dataclass
class RowContext:
    rows:              List[Dict[str, Any]]   # full partition
    variables:         Dict[str, List[int]]   # var -> [row indices]
    current_idx:       int                    # row under test
    current_var:       Optional[str]          # variable being matched
    subsets:           Dict[str, List[str]]   # SUBSET definitions
    defined_variables: Set[str]               # DEFINE variable names
    pattern_variables: List[str]              # all pattern variables
\end{lstlisting}

\noindent
Predicates are evaluated lazily during DFA simulation: a condition
is called only when the DFA attempts the corresponding transition.
For large partitions (over $50\,000$ rows) the engine switches to a
vectorised evaluation path that pre-computes a
\emph{condition matrix}—a Boolean array of shape
$|\text{variables}| \times n$—using NumPy broadcast operations,
reducing condition evaluations from $O(n \cdot |\text{variables}|)$
scalar calls to $O(|\text{variables}|)$ vector operations.

% ................................
\subsection{DFA Simulation}
\label{subsec:meth-exec-dfa-sim}
% ................................

The core of the execution stage is the DFA simulation loop, shown
as Algorithm~\ref{alg:dfa-sim}.  Starting from the DFA's initial
state $q_0$, the engine processes the sorted partition one row at a
time.  At each row $r_i$ it evaluates the \texttt{DEFINE} predicates
for every outgoing transition of the current state, collects the
transitions whose predicates evaluate to \texttt{True}, and selects
the highest-priority one (reflecting greedy or leftmost-match
semantics).  When the DFA reaches an accepting state, the engine
records the current variable bindings as a candidate match and
either terminates (for simple patterns) or continues to seek a
longer greedy match.

\begin{algorithm}[htbp]
\caption{DFA Simulation Loop for a Single Partition}
\label{alg:dfa-sim}
\begin{algorithmic}[1]
\Procedure{SimulatePartition}{$P,\, \mathit{dfa},\, \mathit{skip\_mode},\, \mathit{config}$}
  \State $\mathit{results} \gets [\,]$
  \State $i \gets 0$
  \While{$i < |P|$}
    \State $q \gets q_0$
    \State $\mathit{bindings} \gets \{\}$
    \State $\mathit{best} \gets \bot$
    \State $j \gets i$
    \While{$j < |P|$}
      \State $\mathit{fired} \gets
        \{(v,q') \mid (q \xrightarrow{v} q') \in \mathit{dfa},
          C(v, P[j], \mathit{ctx}_{j,\mathit{bindings}})\}$
      \If{$\mathit{fired} = \varnothing$}
        \State \textbf{break} \Comment{no transition -- end of match attempt}
      \EndIf
      \State $(v^*,q^*) \gets \Call{BestTransition}{\mathit{fired}}$
      \State $\mathit{bindings}[v^*] \gets
          \mathit{bindings}[v^*] \cup \{j\}$
      \State $q \gets q^*$
      \If{$\mathit{dfa.accepting}(q)$}
        \State $\mathit{best} \gets
          (\mathit{start}{=}i,\;\mathit{end}{=}j,\;
           \mathit{bindings}\text{ copy})$
        \If{$\mathit{skip\_mode} = \textsc{ToNextRow}$}
          \State \textbf{break}
        \EndIf
      \EndIf
      \State $j \gets j + 1$
    \EndWhile
    \If{$\mathit{best} \neq \bot$}
      \State $\mathit{results}.\textsc{Append}(\mathit{best})$
      \State $i \gets \Call{NextStart}{i,\;\mathit{best},\;
        \mathit{skip\_mode},\;\mathit{bindings}}$
    \Else
      \State $i \gets i + 1$
    \EndIf
  \EndWhile
  \State \Return $\mathit{results}$
\EndProcedure
\end{algorithmic}
\end{algorithm}

\paragraph{BestTransition selection.}
When multiple transitions fire simultaneously (e.g.\ when a row
satisfies more than one pattern variable's predicate) the engine
applies the following priority ordering, consistent with
SQL:2016 greedy-leftmost semantics:
\begin{enumerate}
  \item Transitions leading to accepting states are preferred over
    non-accepting ones, ensuring the shortest valid prefix is
    always recorded as a candidate match.
  \item Among non-accepting transitions, those whose source and
    target state are the same (self-loops for repeating quantifiers
    such as \texttt{A+}) are preferred over state-advancing
    transitions—implementing greedy extension.
  \item When the pattern contains cross-variable references in
    \texttt{DEFINE} predicates (e.g.\ \texttt{B.price > A.price}),
    the engine further biases toward the self-loop for the
    prerequisite variable (\texttt{A}) before switching to the
    dependent one (\texttt{B}), ensuring $A$'s rows are fully
    consumed before $B$ begins.
  \item Remaining ties are broken by alternation order: the
    leftmost alternative in the original \texttt{PATTERN} string
    receives the lowest numeric priority and is therefore selected
    first.
\end{enumerate}

\paragraph{Variable bindings.}
Throughout the simulation, \texttt{bindings} is a dictionary
$\{ v \mapsto [i_1, i_2, \ldots] \}$ mapping each pattern
variable $v$ to the (sorted) list of partition row indices
assigned to it during the current match attempt.  At the
conclusion of a successful match these lists provide the concrete
row sets over which measure expressions and navigation functions
are evaluated.

Rows belonging to \emph{exclusion} variables (those wrapped in
\texttt{\{- \ldots -\}}) are tracked in a separate
\texttt{excluded\_rows} list: they participate in condition
evaluation and DFA state transitions normally but are omitted from
the measure-evaluation context and from the output rows in
\texttt{ALL ROWS PER MATCH} mode, conforming to the SQL:2016
semantics for pattern exclusions.

% ................................
\subsection{After Match Skip}
\label{subsec:meth-exec-skip}
% ................................

After recording a match the engine must decide where to start the
next match attempt, controlling whether matches can overlap.
SQL:2016 defines four \texttt{AFTER MATCH SKIP} strategies;
Table~\ref{tab:skip-modes} summarises their semantics and the
corresponding next-start-position computation.

\begin{table}[htbp]
  \centering
  \caption{The four \texttt{AFTER MATCH SKIP} modes and the
    resulting next scan position}
  \label{tab:skip-modes}
  \begin{tabular}{lp{9cm}}
    \toprule
    \textbf{Mode} & \textbf{Next-start computation} \\
    \midrule
    \textsc{Past Last Row} (default) &
      $i \leftarrow \mathit{best.end} + 1$. No overlap is possible;
      the scan resumes immediately after the last matched row. \\[4pt]
    \textsc{To Next Row} &
      $i \leftarrow \mathit{best.start} + 1$. Overlapping matches
      are produced because only a single row is skipped regardless
      of match length. \\[4pt]
    \textsc{To First} $V$ &
      $i \leftarrow \min(\mathit{bindings}[V]) + 1$. The scan
      resumes after the \emph{first} row assigned to variable $V$.
      Requires $V$ to appear in the pattern; enforced at compile time. \\[4pt]
    \textsc{To Last} $V$ &
      $i \leftarrow \max(\mathit{bindings}[V]) + 1$. The scan
      resumes after the \emph{last} row assigned to variable $V$.
      May overlap if $V$'s last row precedes the match end. \\
    \bottomrule
  \end{tabular}
\end{table}

\noindent
To prevent infinite loops, the engine enforces a monotonic
progress invariant: the next scan position is always strictly
greater than the current start position.  For
\textsc{To Next Row} the engine additionally detects cycles in
the sequence of recent start positions and breaks out if the same
position is revisited.

% ................................
\subsection{Output Modes}
\label{subsec:meth-exec-output}
% ................................

SQL:2016 specifies two principal output-row multiplicity modes,
each admitting further sub-options
(Table~\ref{tab:rows-per-match}).

\begin{table}[htbp]
  \centering
  \caption{Output modes and their row-count and measure-semantics
    behaviour}
  \label{tab:rows-per-match}
  \begin{tabular}{llp{7cm}}
    \toprule
    \textbf{Mode} & \textbf{Rows out} & \textbf{Measure evaluation context} \\
    \midrule
    \textsc{One Row Per Match}
      & 1 per match
      & \textsc{Final}: measures see the complete, finalised
        variable bindings.  Navigation functions
        (\texttt{FIRST}, \texttt{LAST}) operate over the full
        matched span. \\[6pt]
    \textsc{All Rows Per Match}
      & 1 per matched row
      & Default semantics per function: \texttt{LAST} uses
        \textsc{Running}; \texttt{FIRST}/\texttt{PREV}/\texttt{NEXT}
        use \textsc{Final}; aggregates
        (\texttt{COUNT}, \texttt{SUM}) use \textsc{Running}.
        Explicit \texttt{RUNNING}/\texttt{FINAL} keywords override
        the default. \\[6pt]
    \textsc{All Rows \ldots With Unmatched}
      & 1 per every input row
      & Matched rows carry measure values; unmatched rows emit
        \texttt{NULL} for all measure columns. \\
    \bottomrule
  \end{tabular}
\end{table}

\paragraph{One Row Per Match.}
For each recorded match the engine constructs a single output
dictionary.  The \texttt{PARTITION BY} and \texttt{ORDER BY}
column values are copied from the first matched row; measure
expressions are evaluated once with \textsc{Final} semantics
(all variable bindings are visible).  The resulting dictionary is
appended to the results list and later assembled into the output
DataFrame.

\paragraph{All Rows Per Match.}
For each recorded match the engine iterates over every row
belonging to the match span (from \texttt{best.start} to
\texttt{best.end} inclusive).  For each row $r_j$ in the span
the engine:
\begin{enumerate}
  \item assigns $r_j$ a \texttt{CLASSIFIER} column value equal to
    the name of the pattern variable to which row $j$ was assigned
    (or \texttt{NULL} if the row was excluded);
  \item evaluates each measure expression with the semantics
    determined by Table~\ref{tab:rows-per-match}: \textsc{Running}
    semantics limits variable bindings to rows $\leq j$,
    while \textsc{Final} semantics always uses the complete
    bindings; and
  \item appends the resulting row dictionary to the results list.
\end{enumerate}
Rows covered by the match span but assigned to an exclusion
variable are handled according to the SQL:2016 rule: they appear
in the output with a \texttt{NULL} \texttt{CLASSIFIER} value and
\texttt{NULL} measure columns.

% ................................
\subsection{Measure Evaluation}
\label{subsec:meth-exec-measures}
% ................................

Measure expressions are evaluated by the
\texttt{MeasureEvaluator} class in
\texttt{src/matcher/measure\_evaluator.py}.  Each declared
measure is compiled into a callable similar to the condition
callables described in Section~\ref{subsec:meth-exec-classify}.
The evaluator dispatches on the expression type:

\begin{description}
  \item[Navigation functions.]
    \texttt{FIRST($v$.$c$)}, \texttt{LAST($v$.$c$)},
    \texttt{PREV($c$, $k$)}, and \texttt{NEXT($c$, $k$)}
    resolve to direct index lookups into the partition row list:
    \begin{itemize}
      \item $\texttt{FIRST}(v.c) \to
        P[\mathit{bindings}[v][0]][c]$
      \item $\texttt{LAST}(v.c) \to
        P[\mathit{bindings}[v][-1]][c]$
      \item $\texttt{PREV}(c,\,k) \to
        P[j - k][c]$ (where $j$ is the current output row index)
      \item $\texttt{NEXT}(c,\,k) \to
        P[j + k][c]$
    \end{itemize}
    All accesses are bounds-checked: out-of-partition references
    return \texttt{NULL} per SQL:2016.

  \item[Aggregate functions.]
    \texttt{COUNT}, \texttt{SUM}, and \texttt{AVG} are implemented
    as running or final aggregates over the row index sets stored
    in \texttt{bindings}.  Under \textsc{Running} semantics the
    aggregate is computed over
    $\{ i \in \mathit{bindings}[v] \mid i \leq j \}$;
    under \textsc{Final} semantics it uses all rows in
    $\mathit{bindings}[v]$.

  \item[Pattern-variable column references.]
    An expression of the form $v.c$ (without a function wrapper)
    resolves to the column $c$ of the \emph{first} row assigned to
    $v$.  In \textsc{Running} mode the reference is restricted to
    assigned rows up to the current output position.

  \item[Arithmetic and logical expressions.]
    Compound expressions (e.g.\ \texttt{LAST(A.price) -
    FIRST(A.price)}) are evaluated by recursively evaluating
    sub-expressions and applying the indicated operator.
\end{description}

\paragraph{RUNNING vs.\ FINAL semantics.}
The distinction between the two semantics is central to
\texttt{ALL ROWS PER MATCH} output.  Under \textsc{Running}
semantics a measure evaluated for output row $j$ sees only rows
$r_0, \ldots, r_j$—the match is \emph{in progress} and future
rows are unknown.  Under \textsc{Final} semantics the full
match is always available regardless of which row is being output.

Consider a \texttt{SUM(A.price)} measure over a match covering
rows 3, 5, and 7.  In \textsc{Running} mode the output for each
row is a cumulative partial sum:
\begin{center}
\begin{tabular}{cl}
  Output row & \textsc{Running} \texttt{SUM} \\
  \hline
  3 & $P[3].\mathit{price}$ \\
  5 & $P[3].\mathit{price} + P[5].\mathit{price}$ \\
  7 & $P[3].\mathit{price} + P[5].\mathit{price} + P[7].\mathit{price}$ \\
\end{tabular}
\end{center}
whereas \textsc{Final} mode emits the same total on all three
output rows.

% ................................
\subsection{Worked End-to-End Example}
\label{subsec:meth-exec-example}
% ................................

To illustrate the full execution pipeline consider the following
query applied to a stock-price DataFrame with columns
\texttt{symbol}, \texttt{date}, and \texttt{price}:

\begin{lstlisting}[language=SQL, caption={Example query: detecting
  a V-shape price pattern}, label={lst:exec-example-query}]
SELECT *
FROM prices
MATCH_RECOGNIZE (
  PARTITION BY symbol
  ORDER BY date
  MEASURES
    FIRST(A.price) AS start_price,
    LAST(B.price)  AS bottom_price,
    LAST(C.price)  AS end_price
  ONE ROW PER MATCH
  AFTER MATCH SKIP PAST LAST ROW
  PATTERN (A+ B+ C+)
  DEFINE
    A AS A.price >= PREV(A.price, 1),
    B AS B.price <  PREV(B.price, 1),
    C AS C.price >  PREV(C.price, 1)
)
\end{lstlisting}

\noindent
Suppose the partition for \texttt{symbol = 'AAPL'} contains
five rows sorted by \texttt{date} with prices
$100, 102, 98, 95, 103$.  Execution proceeds as follows.

\vspace{0.5em}
\noindent
\textbf{Step 1 — Partition and Order.}
\texttt{groupby('symbol')} produces a single-symbol sub-frame;
\texttt{sort\_values('date')} confirms the price sequence
$[100, 102, 98, 95, 103]$.

\vspace{0.5em}
\noindent
\textbf{Step 2 — DFA construction.}
The pattern \texttt{A+ B+ C+} compiles to an NFA with 8 states via
Thompson's construction and then to a DFA with 4 states via subset
construction (Section~\ref{subsec:meth-dfa}).

\vspace{0.5em}
\noindent
\textbf{Step 3 — Simulation from $i=0$.}

\begin{itemize}[leftmargin=2em]
  \item Row 0 ($p=100$): \texttt{PREV(A.price)} is undefined at the
    first row; the engine treats this as trivially satisfying
    $A$'s predicate.  Transition $q_0 \xrightarrow{A} q_1$ fires;
    $\mathit{bindings}[A] = [0]$.  State $q_1$ is not accepting.
  \item Row 1 ($p=102 \geq 100$): \texttt{A} predicate satisfied;
    self-loop $q_1 \xrightarrow{A} q_1$.  $\mathit{bindings}[A] = [0,1]$.
  \item Row 2 ($p=98 < 102$): \texttt{A} fails, \texttt{B} passes
    ($p < \mathit{prev}$).  Transition $q_1 \xrightarrow{B} q_2$.
    $\mathit{bindings}[B] = [2]$.
  \item Row 3 ($p=95 < 98$): \texttt{B} satisfied; self-loop
    $q_2 \xrightarrow{B} q_2$.  $\mathit{bindings}[B] = [2,3]$.
  \item Row 4 ($p=103 > 95$): \texttt{C} satisfied; transition
    $q_2 \xrightarrow{C} q_3$.  $q_3$ is accepting.
    $\mathit{best} = (\mathit{start}{=}0,\; \mathit{end}{=}4,\;
    \{A{:}[0,1],\; B{:}[2,3],\; C{:}[4]\})$.
\end{itemize}

\vspace{0.5em}
\noindent
\textbf{Step 4 — After Match Skip.}
Mode is \textsc{Past Last Row}: next scan starts at $i = 5$, which
equals $|P|$, so the simulation terminates.

\vspace{0.5em}
\noindent
\textbf{Step 5 — Measure evaluation (ONE ROW PER MATCH, FINAL).}
\begin{align*}
  \texttt{start\_price}  &= \texttt{FIRST}(A.\mathit{price})
    = P[0][\mathit{price}] = 100 \\
  \texttt{bottom\_price} &= \texttt{LAST}(B.\mathit{price})
    = P[3][\mathit{price}] = 95 \\
  \texttt{end\_price}    &= \texttt{LAST}(C.\mathit{price})
    = P[4][\mathit{price}] = 103
\end{align*}

\vspace{0.5em}
\noindent
\textbf{Step 6 — Output assembly.}
A single output row is emitted:
\texttt{(symbol='AAPL', start\_price=100, bottom\_price=95,
end\_price=103)}.

This example demonstrates how all execution components--partition
ordering, DFA simulation with greedy quantifiers, variable binding
accumulation, after-match skip, and measure evaluation—interact to
produce a single SQL:2016-compliant output row from a five-row
price sequence.


% ----------------------------------------------------------------
\section{Implementation and Evaluation Metrics}
\label{sec:meth-impl}
% ----------------------------------------------------------------

The implementation provides a \texttt{MATCH\_RECOGNIZE} engine for Pandas
DataFrames following the SQL:2016 standard for pattern matching in
sequential data.  The implementation has four core components: an
ANTLR4-based SQL parser that transforms \texttt{MATCH\_RECOGNIZE}
queries into abstract syntax trees; a finite-automata engine that
converts pattern expressions into optimised NFAs and DFAs using
Thompson's construction; a pattern-matching executor that processes
DataFrames with support for partitioning and ordering; and a measure
evaluator that handles navigation functions (\texttt{PREV},
\texttt{NEXT}, \texttt{FIRST}, \texttt{LAST}) and aggregations with
both running and final semantics.  The system supports advanced
pattern features including anchors, quantifiers, alternation,
permuted patterns, and exclusion syntax.

% ................................
\subsection{ANTLR4 Grammar and AST Construction}
\label{subsec:impl-antlr4}
% ................................

The system uses a customised SQL grammar extended from Trino's SQL
that fully supports the SQL:2016 \texttt{MATCH\_RECOGNIZE} clause.
This grammar was defined using ANTLR4, covering complex patterns,
nested patterns, alternation, quantifiers, grouping, and
exclusions~\cite{antlr_trino_sql}.  The parser's output is a
detailed parse tree that captures the syntactic structure of each
SQL query~\cite{parr2013definitive,tarazona2021antlr}.

% ................................
\subsection{Finite Automata}
\label{subsec:impl-automata}
% ................................

\paragraph{Pattern Tokenisation.}
Pattern expressions are tokenised using a recursive-descent
tokenisation algorithm that supports anchors, grouping,
quantifiers, and nested constructs.

\paragraph{NFA Construction.}
Thompson's construction algorithm is used to translate pattern
tokens into a Non-deterministic Finite Automaton (NFA).  Each
state captures transitions, variable bindings, anchors, and
metadata for advanced features.  Quantifiers (\texttt{*},
\texttt{+}, \texttt{?}, \texttt{\{$n$,$m$\}}) are directly
mapped to sub-automata~\cite{georgiou2025streaming}.

\paragraph{DFA Construction and Optimisation.}
The NFA is systematically converted into a Deterministic Finite
Automaton (DFA) using subset construction, drastically reducing
the number of active states during execution.  Optimisation steps
including state minimisation, transition indexing, and merging of
equivalent states yield a compact, performant automaton.  The
resulting DFA supports efficient row-by-row evaluation while
preserving pattern semantics, priorities, and custom conditions
for accurate pattern
semantics~\cite{e104-d_3_381,10.1145/986537.986588,georgiou2025streaming}.

% ................................
\subsection{Execution}
\label{subsec:impl-execution}
% ................................

The matching engine applies the DFA sequentially over each partition
of the DataFrame with specialised handling for anchors, skip modes,
and edge cases.  The executor coordinates the entire
pattern-matching process, integrating parsing, automata
construction, match identification, and result formatting.

\paragraph{Condition and Measure Evaluation.}
A dedicated evaluation subsystem translates SQL pattern conditions
and measures into executable Python code that supports navigation
functions and aggregations with SQL:2016-compliant \texttt{NULL}
semantics for both running and final calculations.

% ----------------------------------------------------------------
\section{Performance Optimization Techniques}
\label{sec:bg-performance}
% ----------------------------------------------------------------

% ................................
\subsection{Parallel Partition Processing}
\label{subsec:bg-parallel}
% ................................

Because \texttt{MATCH\_RECOGNIZE} evaluates each
\texttt{PARTITION BY} group independently, partitions provide a
natural unit of parallel work.  The implementation includes a
\texttt{ParallelExecutionManager} in
\texttt{src/utils/performance\_optimizer.py}, which can dispatch
partition-level work items to a thread or process pool when
parallel processing is enabled and the input size is large enough
to justify the scheduling overhead.  If the workload is too small,
system resources are constrained, or parallel execution is disabled
by configuration, the engine falls back to sequential partition
processing.

This optimisation does not change match semantics: each partition
is still ordered, classified, and matched independently, and the
final output is assembled after all partition-level results have
been collected.  Parallelism therefore improves wall-clock time for
multi-partition workloads while preserving the same logical result
as single-threaded execution.

% ................................
\subsection{LRU Pattern Caching}
\label{subsec:bg-caching}
% ................................

Pattern compilation---tokenisation $\to$ NFA construction $\to$ DFA
conversion $\to$ minimisation---is computationally expensive,
especially for patterns with quantifiers and \texttt{PERMUTE}
constructs.  The implementation maintains a thread-safe
\emph{Least Recently Used} (LRU) cache that maps a hash of the
pattern string and \texttt{DEFINE} predicates to the compiled
pattern (the full NFA and DFA pipeline result).
Key configuration parameters are listed in
Table~\ref{tab:bg-cache-config}.

\begin{table}[htbp]
  \centering
  \caption{Pattern cache configuration defaults
    (\texttt{src/config/production\_config.py})}
  \label{tab:bg-cache-config}
  \begin{tabular}{lrl}
    \toprule
    \textbf{Parameter}         & \textbf{Default} & \textbf{Description} \\
    \midrule
    \texttt{cache\_size\_limit}        & 50\,000 entries  & Maximum cached patterns \\
    \texttt{cache\_memory\_limit\_mb}  & 2\,048 MB        & Maximum cache memory \\
    \texttt{cache\_ttl\_seconds}       & 3\,600 s         & Entry time-to-live \\
    \texttt{cache\_clear\_threshold\_mb} & 1\,600 MB      & Proactive eviction threshold \\
    \bottomrule
  \end{tabular}
\end{table}

% ................................
\subsection{Memory Management}
\label{subsec:bg-memory}
% ................................

Large-scale matching over millions of rows requires careful memory
management.  The implementation processes one partition at a time,
keeping only the current partition's data in memory at any point.
Memory usage is monitored via \texttt{psutil}; garbage collection
is triggered proactively when consumption approaches
user-configured limits (\texttt{cache\_clear\_threshold\_mb}).
An object-pool layer (\texttt{src/utils/memory\_management.py})
reuses frequently allocated objects across successive match
operations, reducing allocator pressure on long-running queries.

% ................................
\subsection{Vectorised Condition Evaluation}
\label{subsec:bg-vectorised}
% ................................

For large partitions the engine avoids repeated Python function-call
overhead by pre-computing a \emph{condition matrix}: a dictionary
mapping each pattern variable $v$ to a NumPy Boolean array of
length $n$ (the partition size), where entry $i$ is \texttt{True}
if row $i$ satisfies the \texttt{DEFINE} predicate for $v$.

The pre-computation pass in
\texttt{EnhancedMatcher.\_vectorize\_condition\_evaluation()}
evaluates each variable's predicate function once per row and
stores the results as a compact \texttt{numpy.ndarray(dtype=bool)}.
When the Polars library is available it is used to parallelise this
pass; otherwise a plain Python loop with NumPy storage provides the
same result.  Implicit variables (present in the pattern but absent
from \texttt{DEFINE}) are assigned an all-\texttt{True} array,
matching any row unconditionally.

Once the condition matrix is built, the DFA simulation loop
replaces every per-row predicate call with a single array index
look-up $M[v][j]$, reducing the inner-loop cost from an interpreted
function invocation to a direct memory access.

For partitions below $5\,000$ rows an enhanced
condition-evaluation cache is used instead; for very large
partitions ($>50\,000$ rows) a sample-based pre-warming step
warms internal caches before the full evaluation pass begins.
Both thresholds are hardcoded in
\texttt{src/matcher/matcher.py}.

Table~\ref{tab:perf-summary} summarises all performance
optimisations and the conditions under which each is applied.

\begin{table}[htbp]
  \centering
  \caption{Performance optimisations and their activation
    conditions}
  \label{tab:perf-summary}
  \begin{tabular}{p{3.8cm}p{4.2cm}p{4.4cm}}
    \toprule
    \textbf{Technique} & \textbf{Activation condition}
      & \textbf{Primary benefit} \\
    \midrule
    LRU pattern cache
      & Always on
      & Eliminates redundant NFA/DFA construction across repeated
        queries \\[4pt]
    Parallel partition execution
      & Enabled by configuration and sufficient partition workload
      & Reduces wall-clock time for independent multi-partition
        queries \\[4pt]
    Condition matrix (vectorised look-up)
      & Partition size $> 5\,000$ rows
      & Replaces per-row predicate calls with array index
        look-ups inside the DFA loop \\[4pt]
    One-partition-at-a-time streaming
      & Always on
      & Bounds peak memory to one partition regardless of total
        dataset size \\[4pt]
    Proactive GC / memory guard
      & Process RSS $\geq$ \texttt{cache\_clear\_threshold\_mb}
      & Prevents OOM errors on very large dataset scans \\
    \bottomrule
  \end{tabular}
\end{table}

% ----------------------------------------------------------------
\section{Chapter Summary}
\label{sec:impl-summary}
% ----------------------------------------------------------------

This chapter has presented the implementation and optimisation of the
\texttt{MATCH\_RECOGNIZE} engine.  It explained how row patterns are
tokenised, compiled into NFAs, converted into DFAs, and executed over
ordered Pandas partitions.  It also described the implementation of
after-match skip behaviour, output modes, measure evaluation, and the
runtime components that connect the parser and logical plan to the
matching engine.

The chapter also described the optimisation layer: LRU pattern
caching to avoid repeated compilation, optional parallel
partition processing, partition-wise memory management, and
vectorised condition evaluation to reduce per-row predicate
overhead.  These optimisations are orthogonal to the core matching
semantics and can be evaluated independently in the experimental
chapter.

Chapter~\ref{chap:evaluation} evaluates these design and
implementation decisions empirically, measuring correctness,
throughput, memory behaviour, and performance across the selected
pattern families and dataset sizes.
